This Section discusses estimations of the uncertainties on the
  measured widths \iftoggle{IncludeYields}{and Yields}{} of the dimuon-pairs.
The widths are characterized by the standard deviation of the $
  Y^{\mathrm{corr}} (\Dphi)$ in Eq.\eqref{eq:signal_func}.
Systematics are also evaluated for the $\Gamma$ parameter
  in Eq.~\eqref{eq:signal_func}.


The following sources of systematic uncertainties on the measurements 
  are considered:
\begin{enumerate}
\iftoggle{IncludeYields}
{
  \item Luminosity and \TAA\ uncertainties
}{}
\item Choice of muon selection: Tight (default) vs Medium.
\item Effect of uncertainties on reconstruction efficiencies.
\item Dependence on the \Deta\ cut.
\item Dependence on the Fit model (excluded for the $\Gamma$, as that parameter is tied in with the fit-model).
\item Inclusion of higher-order fourier component in fits in~Eq.\eqref{eq:fit_func}.
\item Alternate efficiency for correcting for the \Mmumu\ cuts.
\item Alternate PDFs for correcting for the Drell-Yan pairs.
\item Dependence on choice of the trigger (Check Only).
\end{enumerate}
These items are discussed further below. 
This section mostly discusses the systematics for the \PbPb\ measurements.
The \pp\ systematics are evaluated in Appendix~\ref{sec:systematicsPP}.





%--------------------------------------------------------------------------
\iftoggle{IncludeYields}
{
  \subsection{Luminosity and \TAA\ uncertainties\label{sec:systematics_lumi}}
  The luminosity used in this analysis as an uncertainty of $\pm1.5\%$ for the \PbPb\
    data and 1.6\% for the \pp\ data.
  While the measured yields are not directly affected by the luminosity,
    they are reported as yields corresponding to a luminosity of 1.93~\inb.
  For the \TAA\ scaled per-event yields Eq.~\eqref{eq:TaaYield},
    the $N_{\mathrm{events}}^{\mathrm{PbPb}}$ are obtained
    using the luminosity, so luminosity uncertainties
    are directly relevant.
  This uncertainty is directly propagated as an overall uncertainty 
    on the Yields and \TAA\ scaled per-event yields (or cross-section for \pp).
  Similarly the \TAA\ uncertainties (listed in Table~\ref{tbl:TAA})
    are propagated to the \TAA\ scaled per-event yields.
  These uncertainties do not affect the widths parameters of the fits.
}{}
%--------------------------------------------------------------------------





%--------------------------------------------------------------------------
\subsection{Muon selection\label{sec:systematics_muonselection}}
The default analysis is performed using the ``Tight'' muon working point,
  which gives a purer muon selection compared to the ``Medium''
  working point, but yields lower statistics.
The measurements are repeated using the ``Medium'' working point,
  with the intent of including any systematic variations,
  if observed, in the measurements as a systematic uncertainty.
The variation in the measurements when using the Medium working point 
  are shown in Figure~\ref{fig:SystMuonQuality} for the widths
  and in Figure~\ref{fig:SystMuonQualityTau} for $\Gamma$.
\iftoggle{IncludeYields}
{
  They are shown in Figure~\ref{fig:SystMuonQuality2} for the Yields.
}{}
This comparison covers effects of possible contamination of the
  measurement from fakes and from misidentified muons -- which
  are increased in the Medium working point.
The systematic uncertainty is smoothed using the fits
  shown in the lower panels of 
  \iftoggle{IncludeYields}
  {
    Figures~\ref{fig:SystMuonQuality}--\ref{fig:SystMuonQuality2}.
  }
  {
    Figures~\ref{fig:SystMuonQuality} and \ref{fig:SystMuonQualityTau}.
  }
\FuncDescr



\begin{figure}
\centering
\includegraphics[width=1.0\textwidth]%
{figures/AllFigs/can_yields_QualityDep_YieldType0_Tight_pTBar_GT5_EffCor}
\caption{
Top panel: variations in the measured widths of the away-side correlation 
  when using the Tight muon working-point vs. the Medium
  muon working-point.
From left to right, the panels show the comparisons for
  same-sign pairs, opposite-sign pairs and All-pairs.
The lower panels show the variation as a percentage,
  together with a fit function used to smooth 
  the uncertainties.
\FuncDescr\  
%The statistical uncertainties are not plotted in the 
%  lower panels, for clarity.
For determining the fits in the lower panels, each data-point is given equal weight, 
  irrespective of its uncertainty.
}
\label{fig:SystMuonQuality}
\end{figure}


\begin{figure}
\centering
\includegraphics[width=1.0\textwidth]%
{figures/AllFigs/can_yields_QualityDep_YieldType5_Tight_pTBar_GT5_EffCor}
\caption{
Same as Figure~\ref{fig:SystMuonQuality} but for the $\Gamma$.
}
\label{fig:SystMuonQualityTau}
\end{figure}



\iftoggle{IncludeYields}
{
  \begin{figure}
  \centering
  \includegraphics[width=1.0\textwidth]%
  {figures/AllFigs/can_yields_QualityDep_YieldType1_Tight_pTBar_GT5_EffCor}
  \caption{
  Same as Figure~\ref{fig:SystMuonQuality} but for the Yields. 
  }
  \label{fig:SystMuonQuality2}
\end{figure}
}{}
%--------------------------------------------------------------------------





%--------------------------------------------------------------------------
\FloatBarrier
\subsection{Reconstruction and Trigger efficiency uncertainties\label{sec:systematics_efficiencies}}

\begin{figure}
\centering
\includegraphics[width=1.0\textwidth]%
{figures/AllFigs/can_yields_EffDep_YieldType0_Tight_pTBar_GT5_EffCor}
\caption{
Top panel: variations in the measured widths of the away-side correlation 
  when varying the muon reconstrution efficiencies up and down 
  within their allowed uncertainties.
From left to right, the panels show the comparisons for
  same-sign pairs, opposite-sign pairs and All-pairs.
The lower panels show the variation as a percentage,
  together with a fit function used to smooth 
  the uncertainties.
\FuncDescr\ 
The statistical uncertainties are not plotted in the 
  lower panels, for clarity.
For determining the fits in the lower panels, each data-point is given equal weight, 
  irrespective of its uncertainty.
}
\label{fig:SystEfficiency}
\end{figure}

\begin{figure}
\centering
\includegraphics[width=1.0\textwidth]%
{figures/AllFigs/can_yields_EffDep_YieldType5_Tight_pTBar_GT5_EffCor}
\caption{
Same as Figure~\ref{fig:SystEfficiency} but for the $\Gamma$.
}
\label{fig:SystEfficiencyTau}
\end{figure}


\iftoggle{IncludeYields}
{
  \begin{figure}
  \centering
  \includegraphics[width=1.0\textwidth]%
  {figures/AllFigs/can_yields_EffDep_YieldType1_Tight_pTBar_GT5_EffCor}
  \caption{
  Same as Figure~\ref{fig:SystEfficiency} but for the Yields. 
  }
  \label{fig:SystEfficiency2}
  \end{figure}
}{}

The effects of reconstruction efficiencies are
  accounted for by  applying corresponding weights
  -- inverse of the efficiencies --  as a function of \pt\
  and $q*\eta$ to the muons used in the analysis.
The efficiency measurements were detailed
  in Section~\ref{sec:correction}.
These corrections apply a higher weight to reconstructed muons
  to account for muons that were not reconstructed.
Uncertainties in the efficiencies thus result in uncertainties
  in the measured \Dphi\ distributions and thus in the measured 
  widths/yields of the away-side correlation.
The effect of the reconstruction efficiency uncertainties on the
  meaasurements is determined by varying the muon reconstrution 
  efficiencies up and down within their allowed uncertainties, 
  and repeating the analysis.
Any differences observed in the measurements in included 
  as a systematic uncertainty. 
Figures~\ref{fig:SystEfficiency} and~\ref{fig:SystEfficiencyTau} 
  show the variations in the measured 
  widths and $\Gamma$, respectively, when following the above procedure.
The efficiency systematic mostly leads to a variation in the overall normalization 
  of the pair-yields, but does not siginficantly affect the stape of the distributions.
Consequently the effect of this uncertainty on the widths is quite small.
\iftoggle{IncludeYields}
{
  Figure~\ref{fig:SystEfficiency2} shows the variations in the yields.
}
{}
The systematic uncertainty is obtained using the fits shown in the lower panels of 
\iftoggle{IncludeYields}{
  Figures~\ref{fig:SystEfficiency}--\ref{fig:SystEfficiency2}.
  In addition to these uncertainties for the muon efficiency, 
    an additional uncertainty of 1\% is included, which accounts 
    for the muon ID reconstruction efficiency systematics~\cite{HION-2019-11}.
  \FuncDescr\ 
  Effects from the the trigger efficiency are also integrated into the 
    systematic uncertainties.
  The trigger efficiency uncertainty is included based on the difference between
    the change in the efficiency between most central and most-peripheral
    bins for medium and tight muons:
  \begin{eqnarray}
  \delta\epsilon=\frac{1}{2}(\Delta\epsilon_{\mathrm{medium}}-\Delta\epsilon_{\mathrm{tight}}),\nonumber
  \end{eqnarray} 
  with:
  \begin{eqnarray}
  \Delta\epsilon_{\mathrm{medium}}=\epsilon^{\mathrm{0-10\%}}-\epsilon^{\mathrm{60--100\%}}.\nonumber
  \end{eqnarray}
  These uncertainties are $\sim3\%$ leading to $\sim$6\% uncertainty for the pairs.
}
{
  Figures~\ref{fig:SystEfficiency} and~\ref{fig:SystEfficiencyTau}.
}
As an additional check (not included in systematic uncertainties),
  Figures~\ref{fig:SystEfficiency2} and~\ref{fig:SystEfficiency2Tau}
  show the variation in the width and $\Gamma$ when not applying 
  any efficiency corrections at all.
The variation is quite minimal and consistent with the nominal measurements
  within statistical uncertainties.



\begin{figure}
\centering
\includegraphics[width=1.0\textwidth]%
{figures/AllFigs/can_yields_EffDep2_YieldType0_Tight_pTBar_GT5_EffCor}
\caption{
Top panel: variations in the measured widths of the away-side correlation 
  when not applying the trigger and reconstruction efficiency corrections.
From left to right, the panels show the comparisons for
  same-sign pairs, opposite-sign pairs and All-pairs.
The lower panels show the variation as a percentage,
  together with a fit function used to smooth 
  the uncertainties.
\FuncDescr\ 
}
\label{fig:SystEfficiency2}
\end{figure}

\begin{figure}
\centering
\includegraphics[width=1.0\textwidth]%
{figures/AllFigs/can_yields_EffDep2_YieldType5_Tight_pTBar_GT5_EffCor}
\caption{
Same as Figure~\ref{fig:SystEfficiency2} but for the $\Gamma$.
}
\label{fig:SystEfficiency2Tau}
\end{figure}
%--------------------------------------------------------------------------







%--------------------------------------------------------------------------
\FloatBarrier
\subsection{Dependence in the \Deta\ cut \label{sec:systematics_deta}}
As discussed before in Section~\ref{sec:muon_pair_selections}, 
  a requirement of \Deta>0.8 is imposed on the two muons in the pair,
  in order to remove the large near-side jet peak from the measurements.
A systematic uncertainty is included to account for possible
  biases that this requirement introduces.
The measurement is repeated with a stronger requirement of \Deta>0.9. 
Figure~\ref{fig:SystDeta} compares the changes in the measured widths, 
  when performing this variation.
Figure~\ref{fig:SystDetaTau} shows the corresponding plots for the $\Gamma$.
\iftoggle{IncludeYields}
{
  Figure~\ref{fig:SystDeta2} shows the corresponding variations for the yields.
  For the yields, the increase of the \Deta\ gap from 0.8 to 0.9 results in a 
    geometric reduction in the pair-acceptance.
  The geometric reduction is $(4.0^2 -3.9^2)/4.0^2\sim5\%$,
    where $4.0^2$ is the original pair-acceptance (with \Deta>0.8)
    and  $3.9^2$ is the reduced pair acceptance (with \Deta>0.9).  
  When considering the effect that the \Deta\ separation increase has on the yields,
    this trivial factor of 5\% is taken out.
}
The systematic uncertainty is obtained using the fits shown in the lower panels of 
  \iftoggle{IncludeYields}{
    Figures~\ref{fig:SystDeta}--\ref{fig:SystDeta2}.
  }
  {
    Figures~\ref{fig:SystDeta} and~\ref{fig:SystDetaTau}.
  }
\FuncDescr\ 

\begin{figure} 
\centering
\includegraphics[width=1.0\textwidth]%
{figures/AllFigs/can_yields_DetaDep_YieldType0_Tight_pTBar_GT5_EffCor}
\caption{
Top panel: variations in the measured widths of the away-side correlation 
  when applying a $\Deta$>0.8 (default) requirement on the 
  muons in the pairs, vs. when applying a $\Deta$>0.9 requirement.
From left to right, the panels show the comparisons for
  same-sign pairs, opposite-sign pairs and All-pairs.
The lower panels show the variation as a percentage,
  together with a fit function used to smooth 
  the uncertainties.
\FuncDescr\ 
%The statistical uncertainties are not plotted in the 
%  lower panels, for clarity.
For determining the fits in the lower panels, each data-point is given equal weight, 
  irrespective of its uncertainty.
}
\label{fig:SystDeta}
\end{figure}

\begin{figure}
\centering
\includegraphics[width=1.0\textwidth]%
{figures/AllFigs/can_yields_DetaDep_YieldType5_Tight_pTBar_GT5_EffCor}
\caption{
Same as Figure~\ref{fig:SystDeta} but for the $\Gamma$.
}
\label{fig:SystDetaTau}
\end{figure}


\iftoggle{IncludeYields}{
  \begin{figure}
  \centering
  \includegraphics[width=1.0\textwidth]%
  {figures/AllFigs/can_yields_DetaDep_YieldType1_Tight_pTBar_GT5_EffCor}
  \caption{
  Same as Figure~\ref{fig:SystDeta} but for the Yields.
  The fit function in this case is taken to be a linear function
    over the 10--80\% centrality range. 
  }
  \label{fig:SystDeta2}
  \end{figure}
}
{}
%--------------------------------------------------------------------------




%--------------------------------------------------------------------------
\FloatBarrier
\subsection{Dependence on the Fit model\label{sec:systematics_fitmodel}}
The analysis used Eq.~\ref{eq:fit_func}--~\ref{eq:signal_func} to extract
  the number of correlated pairs and the width of the correlation.
Here an alternate model-independent method is used to evaluate the correlated 
  pairs and the correlation width.
For the alternate model, the measured correlation  for 
  $-\pi/2<\Dphi<\pi/2$ is shifted by $\pi$ and subtracted from
  the correlation measured in the range $\pi/2<\Dphi<3\pi/2$.
After this subtraction the 
  \iftoggle{IncludeYields}{yields and widths}{widths}
  on the away-side ($\pi/2<\Dphi<3\pi/2$) are directly evaluated from 
  the subtracted histograms.
The differences for results obtained with this \textit{Subtracted} procedure and the 
  Nominal one (Eq.~\ref{eq:fit_func}--~\ref{eq:signal_func})
  are quite significant.
This is because this new procedure implicitly assumes that the 
  correlation becomes zero for $-\pi/2<\Dphi<\pi/2$,
  while the $Y^{\mathrm{corr}}(\Dphi)$ in the nominal procedure
  has a significant tail that extends below $\pi/2$.
In order to do a more consistent comparison, the nominal-results
  are re-evaluated with a similar subtraction applied.
That is, with the $Y^{\mathrm{corr}}(\Dphi<\pi/2)$, shifted and subtracted from
  $Y^{\mathrm{corr}}(\Dphi>\pi/2)$, 
  \iftoggle{IncludeYields}{and the integral and rms-widths}{and the rms-widths}
  computed with after this subtraction.
The fractional difference between results obtained from these 
  two procedures is taken as a systematic uncertainty.
The systematic uncertainty obtained from this procedure
  is summarized in 
  \iftoggle{IncludeYields}{Figures~\ref{fig:SystModel} and ~\ref{fig:SystModel2}}{Figure~\ref{fig:SystModel}}.
This uncertainty is excluded for the $\Gamma$, as that parameter is 
  tied in with the fit-model (i.e. the parameter is part of the fit model).


\begin{figure}
\centering
\includegraphics[width=1.0\textwidth]%
{figures/AllFigs/can_yields_ModelDep_YieldType0_Tight_pTBar_GT5_EffCor}
\caption{
Top panel: variations in the measured widths of the away-side correlation 
  when 
  1) using the nominal fitting procedure (Eq.~\ref{eq:fit_func}--~\ref{eq:signal_func})
     but with the $Y^{\mathrm{corr}}(\Dphi<\pi/2)$, shifted and subtracted from
     $Y^{\mathrm{corr}}(\Dphi>\pi/2)$, and
  2) using the \textit{Subtracted} procedure (see text).
From left to right, the panels show the comparisons for
  same-sign pairs, opposite-sign pairs and All-pairs.
The lower panels show the variation as a percentage,
  together with a fit function used to smooth 
  the uncertainties.
\FuncDescr\ 
%The statistical uncertainties are not plotted in the 
%  lower panels, for clarity.
For determining the fits in the lower panels, each data-point is given equal weight, 
  irrespective of its uncertainty.
}
\label{fig:SystModel}
\end{figure}

\iftoggle{IncludeYields}{
  \begin{figure}
  \centering
  \includegraphics[width=1.0\textwidth]%
  {figures/AllFigs/can_yields_ModelDep_YieldType1_Tight_pTBar_GT5_EffCor}
  \caption{
  Same as Figure~\ref{fig:SystModel} but for the Yields. 
  }
  \label{fig:SystModel2}
  \end{figure}
}{}
%--------------------------------------------------------------------------





%--------------------------------------------------------------------------
\FloatBarrier
\subsection{Inclusion of higher-order fourier component in fits\label{sec:systematics_includeV3}}
By default Eq.\eqref{eq:fit_func} includes only a second-order Fourier term ($\cos(2\Dphi)$).
While the second order flow component is the dominating flow contribution, 
  in general the collective flow results in higher order Fourier terms as well.
To account for this, the fits are repeated including a $\cos(3\Dphi)$ term in Eq.~\eqref{eq:fit_func},
  and the variation in the results is included as a systematic uncertainty.
The effect of the variation on the widths is shown in Figure~\ref{fig:SystIncludeV3}
  and is about a 2\% effect.
Figure~\ref{fig:SystIncludeV3Tau} shows the corresponding plot for the $\Gamma$.

\begin{figure}
\centering
\includegraphics[width=1.0\textwidth]%
{figures/AllFigs/can_yields_CheckV3_YieldType0_Tight_pTBar_GT5_EffCor}
\caption{
Top panel: variations in the measured widths of the away-side correlation 
  when including a third order fourier term in Eq.~\eqref{eq:fit_func}.
From left to right, the panels show the comparisons for
  same-sign pairs, opposite-sign pairs and All-pairs.
The lower panels show the variation as a percentage,
  together with a fit function used to smooth 
  the uncertainties.
}
\label{fig:SystIncludeV3}
\end{figure}

\begin{figure}
\centering
\includegraphics[width=1.0\textwidth]%
{figures/AllFigs/can_yields_CheckV3_YieldType5_Tight_pTBar_GT5_EffCor}
\caption{
Same as Figure~\ref{fig:SystIncludeV3} but for the $\Gamma$.
}
\label{fig:SystIncludeV3Tau}
\end{figure}
%--------------------------------------------------------------------------







%--------------------------------------------------------------------------
\FloatBarrier
\subsection{Alternate efficiency for correcting for the \Mmumu\ cuts\label{sec:systematics_MinvEffType}}
Section~\ref{sec:minv_cut_correction} discussed the two ways of estimating the 
  acceptance effects on the opposite-sign pairs that arise because of 
  the \Mmumu, Acoplanarity and \pt-asymmetry cuts.
These were:
\begin{enumerate}
  \item Default Method: By applying the cuts on the ``same-sign'' pairs, and determining
    what fraction of pairs are kept as a function of \Dphi.
  \item Alternate Method: an event-mixing procedure where single muons from different events
    are used to build the pair-distribution and the effect of the cuts on these
    mixed-event pairs is used to estimate the efficiency.
\end{enumerate}
The variation in the results when using these two estimates 
  is shown in Figures~\ref{fig:SystMinvEffType} (for the widths)
  and~\ref{fig:SystMinvEffTypeTau} (for $\Gamma$)
  and is included as a systematic uncertainty.

\begin{figure}
\centering
\includegraphics[width=1.0\textwidth]%
{figures/AllFigs/can_yields_MinvEffType_YieldType0_Tight_pTBar_GT5_EffCor}
\caption{
Top panel: variations in the measured widths of the away-side correlation 
  when using the two different estimates of the efficiency loss
  introduced by the \Mmumu, Acoplanarity and \pt-asymmetry cuts.
From left to right, the panels show the comparisons for
  same-sign pairs, opposite-sign pairs and All-pairs.
The lower panels show the variation as a percentage,
  together with a fit function used to smooth 
  the uncertainties.
Note: this correction doesnot affect the same-sign pairs
  and thus there is no associated systematic uncertainty
  for the same-sharge pairs.
}
\label{fig:SystMinvEffType}
\end{figure}


\begin{figure}
\centering
\includegraphics[width=1.0\textwidth]%
{figures/AllFigs/can_yields_MinvEffType_YieldType5_Tight_pTBar_GT5_EffCor}
\caption{
Same as Figure~\ref{fig:SystMinvEffType} but for the $\Gamma$.
}
\label{fig:SystMinvEffTypeTau}
\end{figure}
%--------------------------------------------------------------------------









%--------------------------------------------------------------------------
\FloatBarrier
\subsection{Dependence on the PDF used in the Drell-Yan correction\label{sec:systematics_drell_yan}}
As discussed previously in Section~\ref{sec:DY_correction}, 
  the measurements are corrected to account for Drell-Yan pairs.
These corrections are applied for the Opposite-Sign, and consequently 
  the Combined-sign pairs.
The default correction uses the nNNPDF2.0 PDF, with the nCTEQ15 and NNPDF3.0
  used to evaluate systematic uncertainties.
The effect of the variation is shown in Figures~\ref{fig:SystDY}
  and~\ref{fig:SystDYTau} for the widths and for $\Gamma$, respectively.
The results are consistent within statistical uncertainties.
This is treated as a check only (and not included as a systematic uncertainty).

\begin{figure}
\centering
\includegraphics[width=1.0\textwidth]%
{figures/AllFigs/can_yields_PDFDep_YieldType0_Tight_pTBar_GT5_EffCor}
\caption{
Top panel: variations in the measured widths of the away-side correlation 
  when using the different PDFs for the Drell-Yan background estimation. 
From left to right, the panels show the comparisons for
  same-sign pairs (where no corrections are applied),
  opposite-sign pairs and All-pairs.
The lower panels show the variation as a percentage,
  together with a fit function used to smooth 
  the uncertainties. 
The fit is made only over the 10-80\% centrality range 
  to exclude the outlier behavior of the 0-10\%
  centrality bin.
The statistical uncertainties are not plotted in the 
  lower panels, for clarity.
Note: this correction doesnot affect the same-sign pairs
  and thus there is no associated systematic uncertainty
  for the same-sharge pairs.
}
\label{fig:SystDY}
\end{figure}

\begin{figure}
\centering
\includegraphics[width=1.0\textwidth]%
{figures/AllFigs/can_yields_PDFDep_YieldType5_Tight_pTBar_GT5_EffCor}
\caption{
Same as Figure~\ref{fig:SystDY} but for the $\Gamma$.
}
\label{fig:SystDYTau}
\end{figure}
%--------------------------------------------------------------------------






%--------------------------------------------------------------------------
\FloatBarrier
\subsection{Dependence on choice of trigger\label{sec:systematics_trigger}}
As discussed previously in Section~\ref{sec:data_PbPb}, 
  the \PbPb\ measurements are performed using
  the \verb=HLT_mu4_mu4NoL1= and \verb=HLT_2mu4= triggers.
The analysis is repeated when using only the \verb=HLT_2mu4= trigger.
The effect of the variation is shown in Figures~\ref{fig:SystTrigger}
  and~\ref{fig:SystTriggerTau} for the widths and for $\Gamma$, respectively.
The results are consistent within statistical uncertainties.
This is treated as a check only (and not included as a systematic uncertainty).

\begin{figure}
\centering
\includegraphics[width=1.0\textwidth,bb=0bp  140bp 570bp 280bp,clip]%
{figures/AllFigs/can_yields_TrigDep_YieldType0_Tight_pTBar_GT5_EffCor}
\caption{
Top panel: variations in the measured widths of the away-side correlation 
  when using both the HLT\_mu4\_mu4NoL1 and  HLT\_2mu4 triggers,
  vs. when using the HLT\_2mu4 trigger only. 
From left to right, the panels show the comparisons for
  same-sign pairs, opposite-sign pairs and All-pairs.
%The lower panels show the variation as a percentage.
%  together with a fit function used to smooth 
%  the uncertainties. 
%The statistical uncertainties are not plotted in the 
%  lower panels, for clarity.
}
\label{fig:SystTrigger}
\end{figure}

\begin{figure}
\centering
\includegraphics[width=1.0\textwidth,bb=0bp  140bp 570bp 280bp,clip]%
{figures/AllFigs/can_yields_TrigDep_YieldType0_Tight_pTBar_GT5_EffCor}
\caption{
Same as Figure~\ref{fig:SystTrigger} but for the $\Gamma$.
}
\label{fig:SystTriggerTau}
\end{figure}



\iftoggle{IncludeYields}{
  \begin{figure}
  \centering
  \includegraphics[width=1.0\textwidth,bb=0bp  140bp 570bp 280bp,clip]%
  {figures/AllFigs/can_yields_TrigDep_YieldType1_Tight_pTBar_GT5_EffCor}
  \caption{
  Same as Figure~\ref{fig:SystTrigger} but for the Yields. 
  }
  \label{fig:SystTrigger2}
  \end{figure}
}{}
%--------------------------------------------------------------------------







%--------------------------------------------------------------------------
\FloatBarrier
The final systematic uncertainties are summarized in 
  Table~\ref{tab:systematics} for the widths.
  and in Table~\ref{tab:systematics2} for the $\Gamma$.
The total uncertainty is taken as the quadrature sum of individual uncertainties.


\begin{table}[h]
\centering
\begin{tabular}{|c|c|c|c|}
\hline
       & Same-sign & Opposite-sign & All-pairs    \\ 
\hline
Source & \multicolumn{3}{ c|}{Uncertainty [\%]}   \\
\hline
Working point       & 3       & 1.5     & $\sim$0 \\
Efficiency          & $\sim$0 & $\sim$0 & $\sim$0 \\
$\Deta$ cut         & 1       & 2       & 1.5     \\
Fit-Model           & 2.5     & 3.5     & 3.5     \\ %TODO
Higher order $\vn$  & 3       & 1       & 1.5     \\
$\Mmumu$ correction & -       & 1       & 0.5     \\
\hline
\end{tabular}
\caption{
  Summary of systematic uncertainty on the widths ($\sigma$).
  \label{tab:systematics}
}
\end{table}


\begin{table}[h]
\centering
\begin{tabular}{|c|c|c|c|}
\hline
       & Same-sign & Opposite-sign & All-pairs    \\ 
\hline
Source & \multicolumn{3}{ c|}{Uncertainty [\%]}   \\
\hline
Working point       & 8       & 2.0     & 0.5     \\
Efficiency          & $\sim$0 & $\sim$0 & $\sim$0 \\
$\Deta$ cut         & 2       & 1       & 1.5     \\
Higher order $\vn$  & 6.5     & 2.5     & 3.5     \\
$\Mmumu$ correction & -       & 2       & 1.0     \\
\hline
\end{tabular}
\caption{
  Summary of systematic uncertainty on the $\Gamma$.
  \label{tab:systematics2}
}
\end{table}
%--------------------------------------------------------------------------

